% Example resume template. Copy it to resumes/base_resume.tex and replace everything with your own details.
% Rules the bot relies on:
%   - Everything above the Education section (preamble + header) is copied unchanged into every tailored resume.
%   - The body must start with the Education section.
%   - Keep it to one A4 page.
% Compiles with XeLaTeX / LuaLaTeX (Overleaf: Menu > Compiler > XeLaTeX) or with Tectonic.
\documentclass[11pt,a4paper]{article}

\usepackage[a4paper,top=0.4in,bottom=0.4in,left=0.45in,right=0.45in]{geometry}

% -------------------- PACKAGES --------------------
\usepackage{enumitem}
\usepackage{hyperref}
\usepackage{titlesec}
\usepackage{fontspec}
\usepackage{xcolor}
\usepackage{fontawesome5}

% -------------------- FONT --------------------
\setmainfont{Latin Modern Roman}
\renewcommand\normalsize{\fontsize{9.5}{11}\selectfont}
\normalsize

% -------------------- LINKS --------------------
\hypersetup{
    colorlinks=true,
    urlcolor=black,
    pdftitle={Your Name - Software Engineer},
    pdfauthor={Your Name},
    pdfkeywords={Software Engineer, Python, Java, REST API, SQL}
}

% -------------------- SECTION FORMAT --------------------
\titleformat{\section}{\scshape\large\raggedright}{}{0em}{}[\titlerule]
\titlespacing*{\section}{0pt}{5pt}{3pt}

% -------------------- LIST SPACING --------------------
\setlist[itemize]{itemsep=1.5pt,topsep=1.5pt,parsep=0pt,leftmargin=1.3em}

% -------------------- GLOBAL --------------------
\pagenumbering{gobble}
\setlength{\parindent}{0pt}
\setlength{\parskip}{0pt}

\begin{document}

% -------------------- HEADER --------------------
\begin{center}
    {\Huge \scshape Your Name}
\end{center}

\vspace{-2pt}
\noindent
\parbox{\textwidth}{\centering\small
\href{mailto:you@example.com}{\faEnvelope\ \underline{you@example.com}} \,
\faPhone\ \underline{+00 0000000000} \,
\href{https://www.linkedin.com/in/your-profile/}{\faLinkedin\ \underline{LinkedIn}} \,
\href{https://github.com/your-username}{\faGithub\ \underline{GitHub}} \,
\href{https://your-portfolio.example.com/}{\faGlobe\ \underline{Portfolio}}
}

\vspace{4pt}

% -------------------- EDUCATION --------------------
\section{Education}
\textbf{Your University} \hfill \textbf{August 2023 - July 2027}\\
\textit{Bachelor of Technology in Computer Science and Engineering} \hfill \textit{CGPA: 0.00/10}

\textbf{Your School} \hfill \textbf{April 2021 - June 2023}\\
\textit{Class XII} \hfill \textit{00.0\%}

% -------------------- SKILLS --------------------
\section{Technical Skills}
\textbf{Languages:} Python, Java, JavaScript, SQL\\
\textbf{Backend \& Cloud:} FastAPI, Node.js, REST APIs, Docker, AWS\\
\textbf{Databases:} PostgreSQL, MySQL, MongoDB\\
\textbf{Tools:} Git, Linux, CI/CD, Unit Testing

% -------------------- EXPERIENCE --------------------
\section{Experience}
\textbf{Company Name - Software Engineering Intern} \hfill \textbf{Month Year - Month Year}
\begin{itemize}
    \item Built \textbf{what you built} using \textbf{which technologies}, which led to \textbf{a measurable result}.
    \item Improved \textbf{something} by \textbf{a number}, by doing \textbf{what you did}.
\end{itemize}

% -------------------- PROJECTS --------------------
\section{Projects}
\textbf{Project Name} \hfill \href{https://github.com/your-username/project}{\textbf{GitHub}}\\
\textit{Python, FastAPI, PostgreSQL, Docker}
\begin{itemize}[nosep,left=0pt]
    \item What the project does, the part you built, and the result in numbers.
\end{itemize}

% -------------------- ACHIEVEMENTS --------------------
\section{Achievements}
\begin{itemize}
    \item \textbf{Award or rank} - Competition name (Year).
\end{itemize}

% -------------------- INTERESTS --------------------
\section{Interests}
Open Source, Competitive Programming

\end{document}
